% Figure — Modèles d'amplificateur : sortie en courant (Norton) et sortie en tension (Thévenin) (EANA3, chapitre 2).
% D'après les notes externes [notes-2023-c04-06, p. 1] et [tableau-2020-c03, p. 5], redessinée.
% Rendu SVG :
%   latex -output-directory=<dossier temporaire> 02-modeles-quadripoles.tex
%   dvisvgm --no-fonts --exact-bbox -o 02-modeles-quadripoles.svg <dossier temporaire>/02-modeles-quadripoles.dvi
\documentclass[border=4pt]{standalone}
\usepackage[T1]{fontenc}
\usepackage[utf8]{inputenc}
\usepackage{amsmath}
\usepackage{sansmath}
\renewcommand{\familydefault}{\sfdefault}
\usepackage[european, straightvoltages, americancurrents]{circuitikz}
\begin{document}
\sansmath
\begin{circuitikz}[font=\small, >=latex, fleche/.style={->, thick}]
  \foreach \dx/\titre in {0/{(a) sortie en courant (Norton)}, 7.6/{(b) sortie en tension (Thévenin)}} {
    \draw[thick, fill=black!3] (\dx+1.0,-0.4) rectangle (\dx+5.4,2.9);
    \node[anchor=west] at (\dx-0.4,3.35) {\titre};
    \draw (\dx,2.0) node[ocirc]{} -- (\dx+1.4,2.0) to[generic, l=$Z_e$] (\dx+1.4,0) -- (\dx,0) node[ocirc]{};
    \draw[fleche] (\dx+0.1,2.25) -- node[above]{$i_e$} (\dx+0.8,2.25);
    \draw[fleche] (\dx-0.3,0.15) -- node[left]{$v_e$} (\dx-0.3,1.85);
    \draw (\dx+6.0,2.0) node[ocirc]{};  \draw (\dx+6.0,0) node[ocirc]{};
  }
  % (a) Norton
  \draw (3.0,2.0) to[american controlled current source, l=$i_s$] (3.0,0);
  \draw (4.4,2.0) to[generic, l=$Z_s$] (4.4,0);
  \draw (3.0,2.0) -- (6.0,2.0);  \draw (3.0,0) -- (6.0,0);
  \node at (3.0,-1.0) {$i_s = f(v_e \text{ ou } i_e)$};
  % (b) Thévenin
  \draw (10.6,2.0) to[american controlled voltage source, l=$v_s$] (10.6,0);
  \draw (10.6,2.0) to[generic, l=$Z_s$] (12.6,2.0) -- (13.6,2.0);  \draw (10.6,0) -- (13.6,0);
  \node at (10.6,-1.0) {$v_s = f(v_e \text{ ou } i_e)$};
\end{circuitikz}
\end{document}
